\documentclass[12pt]{article}
\usepackage[utf8]{inputenc}
\usepackage[T1]{fontenc}
\usepackage{hyperref}
\title{Central Tarahumara Language Resources for K-12 Curriculum Development}
\author{Compiled by Claude Sonnet 5 (claude-sonnet-5) for Andr\'as Kornai}
\date{}
\begin{document}
\maketitle

\section{Introduction}
Central Tarahumara (Rar\'amuri; ISO 639-3 \texttt{tar}, Glottocode cent2131) is a Uto-Aztecan language of Chihuahua, Mexico, with about 30,000 speakers. This report is based on a single web search. It found strong academic documentation but no school materials, primers, or Central Tarahumara dictionary. Nothing below should be read as evidence that such materials do not exist, only that the search did not surface them.

\section{Primary Educational Resources}
No primers, textbooks, or school programmes in Rar\'amuri were found in the search results. Curriculum developers will need to look outside what was found here, for example to Mexican educational authorities for indigenous education, and to local communities.

\section{Linguistic and Reference Materials}
\begin{itemize}
\item Gabriela Caballero, \emph{A grammar of Choguita Rar\'amuri} (Language Science Press, 2022). It is the first comprehensive grammatical description of Choguita Rar\'amuri, a variety of the language, and is open access. Whether it describes the Central variety as a whole was not established. \url{https://langsci-press.org/catalog/book/324}
\item Open-access copy of the grammar record at Refubium (FU Berlin). \url{https://refubium.fu-berlin.de/handle/fub188/37751?show=full}
\item Open-access copy of the grammar record at OAPEN. \url{https://0-library-oapen-org.catalogue.libraries.london.ac.uk/handle/20.500.12657/60563?show=full}
\item Caballero's documentation guide in \emph{Language Documentation \& Conservation} 11 (2017), which describes the documentary collection of the Choguita Rar\'amuri project. \url{https://scholarspace.manoa.hawaii.edu/items/cfa6f5ec-6c4d-4717-9abe-cb9af05cc50b}
\item Rar\'amuri materials in the ACLS collection at the American Philosophical Society (1931 and 1940 material, including word lists and a draft dictionary, from Jules Henry's work). \url{https://indigenousguide.amphilsoc.org/entry/39960}
\item Tarahumara materials in the Jane M. Rosenthal Papers (1969 manuscript drafts and notes on grammar; the variety is not stated). \url{https://indigenousguide.amphilsoc.org/entry/39272}
\end{itemize}

\section{Religious and Community Resources}
None found in the search results.

\section{Digital and Audio Resources}
The Choguita Rar\'amuri documentary corpus (over 200 hours of recordings, according to the search summary) is reported to be held at the Endangered Languages Archive and at UC Berkeley's Survey of California and Other Indian Languages. I did not see archive URLs in the results, so none are given. The documentation guide above is the entry point.

\section{Curriculum Development Context}
The reference description says the language is stable, is the first language of the whole ethnic community, and is not known to be taught in schools. The open-access grammar and corpus could support teacher and materials development. A practical dictionary, a primer and a community-based orthography review would still need to be located or created.

\end{document}